\section{Rastreador IoT}
Nesta seção será discorrido conceitos iniciais e trabalhos relacionados ao contexto de rastreadores em IoT.

\subsection{Conceitos iniciais}
O trabalho desenvolvido baseia-se em três tecnologias essenciais: o sistema de navegação GPS, sistema inercial IMU e protocolo de comunicação de rede LoRaWAN. A seguir, serão descorridos conceitos iniciais concernente a cada um dos citados.

GPS (Global Position System) é definido em 

Em ~\cite{Ahmad2013}, IMU (Inertial Measurement Unit) é definido como um sistemas de sensores usados em conjunto para medir velocidade, orientação e força gravitacional. Inicialmente, a tecnologia era composta de sensor acelerômetro e giroscópio. O acelerômetro é utilizado para medir velocidade inercial enquanto o giroscópio mede ângulos de rotação. Posteriomente foi implementado ao sistema um sensor magnetrômetro, o qual mede a orientação magnética terrestre e trás mais confibilidade para as leituras do giroscópio.

Tais sensores são manufaturados baseado em MEMS (Micro Eletro Mechanical Systems), que consiste em circuitos integrados atuando juntamente com partes mecânicas móveis. Em ~\cite{Zhanshe2015}, uma revisão acerca do desenvolvimento de pesquisa de giroscópios em MEMS de silício é feita. No artigo, é mencionado o princípio de funcionamento. Em poucas palavras, o sistema é contituído por uma massa de prova incercial que, em condições de trabalho externa, reage ao longo das direções e um elemento sensorial para determinar tais situações.

Neste trabalho será utilizado o módulo MPU-6050, ~\cite{MPUDatasheet}, um dispositivo que combina acelerômetro e giroscópio. Também possui um Digital Motion Processor\texttrademark (DMP) com firmware que atua em tempo de execução e visa garantir os dados ideais, atuando como um SoC (System on Chip). Apresenta barramento I2C dedicado e entrada para magnetrômetro externo.

A comunicação entre o dispositivo e a Internet é providenciada pela rede em LoRaWAN. Em ~\cite{Pastorio2021}, 

\subsection{Trabalhos relacionados}

Nesta subseção serão tratados trabalhos relacionacionados a dispositivos rastreadores IoT. Uma abordagem concernente a LoRaWAN e o dispositivo será discorrida nos próximos três parágrafos. Nos parágrafos posteriores, serão relacionados trabalhos e aplicações onde há a combinação de GPS e IMU  em sistemas e dispositivos.

% LoRa e GPS

Um sistema de rastreio com base em LoRaWAN é proposto em ~\cite{Jaffar2020}. Utilizando o microcontrolador Arduino, Dragino YUN shield e LoRa Shield, o dispositivo proposto encaminha suas coordenadas geográficas periodicamente para um LoRa Gateway que então as registra em um Datalog. Testes da aplicação foram feitos em areas urbanas e suburbanas, constatando resultados mais eficientes em areas suburbanas. Ao fim, é defendido que a plataforma LoRaWAN tem potencial para substituir a insfraestrutura de IoT existente baseada em 802.11 Wireless Standard.

Em ~\cite{Pastorio2021}, uma estratégia híbrida de geolacalização é implementada em uma rede LoRaWAN a fim de trazer mais confiabilidade e disponibilidade em sistemas de rastreio. O estudo realizado propõe que, em caso de falha de GPS, a posição do dispositivo pode ser estimada através de algoritimos de aproximação. Resultados apontam erros entre 151m, no pior caso, e 19m no melhor caso.

No artigo ~\cite{Camargo21}, investigações experimentais são feitas para avaliar quatro sistemas de rastremento, dois dispostivos comerciais e dois dispositivos prototipados, com base em LoRaWAN. Testes foram feitos em velocidade costante e variada e os resultados são expressos em RSSI (received signal strength indication), SNR (signal-to-noise ratio) e SF (spreading factor). É também apresentado os resultados do desenvolvimento e de avaliação da rede LoRaWAN em Toledo acerca da cobertura por toda cidade.

% GPS e IMU

Além da definição de IMU, há em ~\cite{Ahmad2013} uma revisão sobre aplicações desse sistema. Em específico, um tópico sobre a aplicação de IMU em sistemas de navegação. É mencionado trabalhos que combinam o uso de GPS e IMU como ~\cite{Bruggemann2011}, onde este sistema é aplicado em ADM (Aircraft Dynamic Model) para detecção de falha de GPS. Também é mencionado uma abordagem para se estabilizar um drone quadrotor utilizando IMU em ~\cite{Azfar2011}.

Avaliações sobre o uso de GPS em conjunto com IMU, acelerômetro e giroscópio, foram realizados em ~\cite{Listiyana2018}. Segundo os autores, o módulo acelerômetro e giroscópio MPU-6050 apresenta um erro médio de 0.96° em todos eixos de rotação. Em específico, o eixo yaw obteve 1.10° de erro médio. Ademais, o teste de posição estática realizado apresenta um erro médio de 2.66m, com erro mais distante de 5.05m. O teste de rastreio de rota apresenta um caminho condizente com o realizado, embora haja alguns desvios. Em conclusão, é defendido que tal sistema de rastreio funciona bem. %Não é mencionado o modelo do GPS

Ainda além, há em ~\cite{Lopes2023} o teste de eficiência do Filtro de Madgwig em conjunto com o Filtro de Kalman para estimar posição e efetividade utilizando dados de IMU para compensar eventual perda de sinal de GPS. O dispositivo proposto é composto pela placa de protótipo Lilygo TTGO v1.1 com o microcontrolador ESP32 e GPS NEO-6M integrado. O sistema IMU é composto pelo módulo acelerômetro e giroscópio MPU-6050 e o módulo magnetômetro HMC5883L. O sistema proposto equipará a frota de veículos elétricos da polícia local do Espírito Santo - Brasil.

Um sistema de envio de coordenada com base em variação de aceleração é comumente aplicado em dispositivos identificadores de colisão de veículos ~\cite{Abdeen2016, Cismas2017, Azlan2023, Karuna2023}. Em ~\cite{Abdeen2016}, é proposto um dispositivo controlado pelo microcontrolador ARM Cortex-M3 STM32 e possui o módulo acelerômetro e giroscópio MPU-6050, além do módulo GSM SIM900 e demais componentes. O dispositivo transmite a sua coordenada geográfica via solicitação por SMS do usuário ou de forma automática. As informações colhidas pelo sensor IMU são processadas para determinar se houve um acidente. Em caso positivo, a localização do dispositivo é encaminhada ao sistema de monitoramento de forma automática. Em ~\cite{agarwal2020} foi desenvolvido um dispositivo para a prevenção de acidentes e roubo de carros baseado em IoT, este sistema utilizou um Arduino em conjunto com um sensor e GPS para obter a localização do veículo, um sensor ultrassônico para perceber a distância entre o veículo monitorado e outros para evitar colisões, um sensor de vibração para detectar se ocorreu um acidente, e um modem GSM para transmitir alertas para uma rede em caso de emergência.

Outra aplicação comum consiste em dispositivos vestíveis identificadores de quedas ~\cite{Manu2021, Yulastri2023, Hsieh2020}. Em ~\cite{Manu2021}, é proposto um dispositivo vestível equipado com uma placa de protótipo Lilygo TTGO-T-CALL com GPS NEO-6M integrado e módulo acelerômetro e giroscópio MPU- 6050. Com base nas leituras do módulo MPU-6050, é detectado se houve a queda da pessoa que utilizava o dispositivo. Por SMS e ligação à um contato pré configurado, a localização da pessoa acidentada é compartilhada. Este dispositivo é pensado para pessoas geriátricas, pacientes de Alzheimer e pacientes de autismo. Em ~\cite{requejo2023} foi desenvolvido um dispositivo que pode ser vestido como um relógio, para pessoas com leve déficit cognitivo ou em estágios iniciais de demência. O aparelho utiliza um microcontrolador STM32L1, um transdutor LoRa SX1276, sensores de acelerômetro e um sensor BME680 para detecção de gás, mudanças na pressão e temperatura, assim como a umidade do ar.

Dentre outras aplicações de ações conjuntas de GPS e IMU há também a proposta de analisar o comportamento dos motoristas em ~\cite{Ayat2021}. O estudo foi realizado com base em leituras realizadas pelo módulo acelerômetro e giroscópio, como velocidade de frenagem e velocidade de curva. Em comparação, foi utilizado o acelerômetro e giroscópio GY521, o mesmo utilizado no módulo MPU-6050, e o acelerômetro e giroscópio GY86. Semelhantemente, também é proposto em ~\cite{mohammed2023intelligent} um sistema baseado em IoT para análise comportamental do motorista.

Ademais, há também a aplicação de detecção de anomalias em pistas rodoviárias em ~\cite{Tripathi2023, Trindade2015}. Utilizando dados de IMU, machine learning e GPS para localização dos eventos. Em ~\cite{Trindade2015}, um protótipo é implementado com um microcontrolador Arduino Uno, shield GPS/GPRS/GSM e módulo acelerômetro e giroscópio MPU-6050. A proposta é o encaminhamento, via SMS, da posição instantânea do veículo e demais informações acerca das variações de aceleração quando identificado uma perturbação na rodovia. A intenção do trabalho é fornecer dados que auxiliem empresas e órgãos públicos no monitoramento da frota de veículos quanto na manutenção de vias. 

Uma outra aplicação proposta em ~\cite{priya2023} foi de analisar o tráfego e prever situações de congestionamento no trânsito. Para este fim foi utilizado um microcontrolador STM32L072CXT6 com um sensor de GPS Dragino LoraWAN IN865, o sensor de GPS manda a latitude, longitude e o horário de cada coordenada para a rede. Além disso o dispositivo é equipado com sensores de qualidade do ar e clima, pois esses fatores afetam a velocidade dos veículos. Outro artigo que descreve uma solução parecida é em ~\cite{molina2023}, que também utiliza machine learning para a predição da localização de veículos.

A necessidade de utilizar módulo magnetrômetro em IMU é discutida em ~\cite{weygers2023}. É apontado que mesmo em ambientes controlados os campos magnéticos são fortemente afetados por fatores externos, o que levanta a questão de se seria mesmo necessário usar magnetrômetros em unidades de medida inercial. A razão pela qual a substituição do magnetrômetro se torna viável é que a estimativa de inclinação por outros meios está tendo um erro cada vez menor, chegando a uma média de 0.9$^{\circ}$ em estudos recentes.

Em ~\cite{huang2023} é discutida uma solução inovadora para ambientes onde o posicionamento por GPS não pode ser determinado. Para isso é utilizado o módulo de GPS para determinar a posição inicial de um conteiner, essa localização também passa como parâmetro o limite de velocidade na pista que o conteiner se desloca, a partir destas informações e de um sensor magnetrômetro, que mede o ângulo em que a carga se direciona, é possível estimar as próximas posições do objeto fazendo chamadas ao GPS apenas de 15 em 15 minutos. Em ~\cite{thuong2010} é construído um sistema para determinar a posição de um veículo e a distância até um destino definido por uma rota, este sistema utiliza sensores magnetrômetro e giroscópio para determinar a direção em que o veículo se desloca caso haja um desvio.

Uma aplicação para determinar precisamente a localização de veículos em vias de alta velocidade é proposto em ~\cite{wu2023}, para isto é utilizado comunicação LoRa junto com IMU para estimar a posição juntamente com o GPS.

Outra aplicação muito utilizada para rastreadores é com relação à pecuária, muitas pesquisas vêm se concentrando em torno do monitoramento e rastreio de manadas de gado. Em ~\cite{baratta2024} é utilizado um sistema com um microcontrolador STM32L0, um módulo GPS NEO-6M, um sensor acelerômetro MPU-9250 e um transdutor LoRa para o monitoramento da manada. Para estimar a atenuação de caminho dos sinais de monitoramento são utilizados modelos de predição deep learning. Outra aplicação semelhante é proposta em ~\cite{longchamp2023}, onde o objetivo era o menor consumo possível de energia, no artigo foi possível deixar o rastreador ligado a uma bateria por 4 meses, e ainda foi possível extender essa vida útil para 6 meses utilizando energia solar. O artigo em questão utiliza sensores acelerômetro, magnetrômetro, de temperatura e um módulo GNSS, tudo isso integrado com a rede por LoRaWAN para o monitoramento da atividade de vacas. 

Em ~\cite{tugrul2019} é proposta uma solução de Geofencing, onde é possível desenhar uma área de interesse em um mapa e um dispositivo móvel irá fazer avisos caso saia desta área de interesse, isto é interessante para o monitoramento de pessoas frágeis, como é o caso de crianças e de pacientes com doenças neurológicas. O artigo também utiliza técnicas como o filtro Kalman e análise de regressão logística para melhorar a precisão do posicionamento por GPS.

Um smartphone Android é utilizado para fazer o rastreamento de um veículo em ~\cite{verma2023}, o dispositivo manda via GSM as coordenadas e os valores dos sensores acelerômetro, giroscópio e magnetrômetro para um servidor que faz o rastreio em tempo real utilizando um mapa. O trabalho ainda propoe a utilização dos sensores do sistema IMU do aparelho para determinar curvas na rota do veículo e mandar coordenadas de acordo com essas curvas para aumentar a precisão da determinação da rota.


